\section{Uniform lattice geometry}\label{sec:geometry}

In this section a line means its intersection with \(V_n\), and only
lines containing at least two grid points are counted. Write \(L_\ell\)
for its cardinality and \(L_{pq}\) for the length of the line through
distinct \(p,q\). Constants in this section are uniform in the root's
position, including roots approaching a cell boundary as \(n\) changes.

\begin{lemma}[Direction and line counts]\label{lem:lines}
There are at most \(4h\) primitive unoriented directions of height
\(h\). Every grid line of that height has at most \(2n/h\) points.
Consequently, uniformly in \(p\),
\begin{align}
 \sum_{\ell\ni p}L_\ell^2&\le16n^2(1+\log n),\label{eq:line2}\\
 \sum_{\ell\ni p}L_\ell^k&\le
 4(2n)^k\left(1+\frac1{k-2}\right)\qquad(k>2).
 \label{eq:linek}
\end{align}
\end{lemma}
\begin{proof}
If \(v\) is primitive, the integer points on its line are precisely
\(p+\Z v\). Indeed, Bezout coefficients show that a real multiple of
\(v\) with integer coordinates has an integer multiplier. A coordinate
of absolute step \(h\) gives \(L_\ell\le1+(n-1)/h\le2n/h\).
For \(h\ge2\), counting the primitive points on the boundary of the
square of radius \(h\), with opposite signs identified, gives
\(4\varphi(h)\le4h\); height one has four directions. Sum
\(4h(2n/h)^k\), using the harmonic bound when \(k=2\) and an integral
comparison when \(k>2\). The same argument gives
\(\sum_{q\ne p}L_{pq}=\sum_{\ell\ni p}(L_\ell-1)L_\ell
\le\sum L_\ell^2\).
\end{proof}

\begin{lemma}[Two uniform incidence kernels]\label{lem:kernels}
For every fixed \(\xi>0\),
\begin{align}
 \sum_{q\notin\ell(a,b)}L_{aq}L_{bq}&\le C_\xi n^{2+\xi}
       &&(a\ne b),\label{eq:twostar}\\
 \sum_{q\in\ell\cap V_n}L_{pq}&\le C_\xi n^{1+\xi}
       &&(p\notin\ell).
       \label{eq:fixedline}
\end{align}
\end{lemma}
\begin{proof}
The elementary divisor estimate \(\tau(m)\le C_\epsilon m^\epsilon\)
follows by writing \(m=\prod l^{a_l}\). For sufficiently large primes,
\(a+1\le2^a\le l^{\epsilon a}\); the finitely many smaller primes have
finite suprema of \((a+1)/l^{\epsilon a}\). Hence, with
\(\Delta_n=\max_{m\le2n^2}\tau(m)\),
\(\Delta_n(1+\log n)^2\le C_\xi n^\xi\).

For \eqref{eq:twostar}, write \(q-a=e u\), \(q-b=f v\), where
\(e,f>0\) and \(u,v\) are primitive, and put \(w=b-a\).
Then \(eu-fv=w\) and \(f\mid\det(u,w)\). The exclusion of the baseline
ensures \(1\le|\det(u,w)|\le2n^2\), so for fixed \(u\) there are at
most \(\Delta_n\) choices of \(f\). In dyadic height bands
\(H\le h(u)<2H\), \(J\le h(v)<2J\), assume \(H\le J\).
There are \(O(H^2)\) choices of \(u\). The congruence
\(eu\equiv w\pmod f\) confines \(e\) to at most one residue class
modulo \(f\), by Bezout. A coordinate with \(|u_i|\ge H\), together
with \(|v_i|<2J\), confines \(e/f\) to an interval of length at most
\(4J/H\). Thus the number of candidates is
\(O(\Delta_n HJ)\). Their individual weights are at most
\(4n^2/(HJ)\). Interchange \(a,b\) when \(H>J\), then sum the
\(O((1+\log n)^2)\) bands.

For \eqref{eq:fixedline}, the cases of zero or one grid point are
immediate. Otherwise write \(\ell\cap\Z^2=a+\Z v\) with primitive
\(v\), and let \(D=|\det(v,p-a)|\in[1,2n^2]\).
For \(q\in\ell\), put \(g(q)=\gcd(|q_1-p_1|,|q_2-p_2|)\).
Then \(g(q)\mid D\). For each \(e\mid D\), the condition
\(q\equiv p\pmod e\) restricts the line parameter to one residue class.
Within \(\|q-p\|_\infty\le R\), this gives at most
\(2R/(e h(v))+1\) choices. Since
\(g(q)\le\sum_{e\mid g(q)}e\) and relevant \(e\le R\),
\[
 \sum_{\|q-p\|_\infty\le R}g(q)\le3R\tau(D).
\]
In the shell \(R/2<\|q-p\|_\infty\le R\),
\(L_{pq}\le4ng(q)/R\), so the shell contributes at most
\(12n\tau(D)\). Sum dyadic shells and use the divisor estimate.
\end{proof}

\subsection{Parity, chords and a uniform boundary estimate}
Suppose \(F_r\) are constant on a fixed partition into \(R\) open
rectangular cells, vanish on the common boundaries, and lie in
\([0,M]\). Along a primitive line through a root of type \(r\), the
types alternate between \(r\) and \(r+\bar v\).

\begin{lemma}[Lattice-to-chord comparison]\label{lem:chord}
For every root and primitive direction of height \(h\),
\begin{equation}\label{eq:chorderror}
 \left|\sum_{q\in\ell(p,v)\setminus\{p\}}
 F_{\typ(q)}(q/n)
 -\frac n{2h}C_v(F_r+F_{r+\bar v})(p/n)\right|
 \le(4R+1)M.
\end{equation}
Moreover \(A_n=A+O_{R,T}(n^{-1})\).
\end{lemma}
\begin{proof}
Each parity subsequence has major-coordinate spacing \(2h/n\).
The number of points of any shifted one-dimensional lattice in an
open interval differs from interval length divided by spacing by at
most two. Sum over the two types and the \(R\) cells, then remove the
root, costing at most \(M\). A line coinciding with a common boundary
has field identically zero on both sides of \eqref{eq:chorderror}.
For the mass statement, count each coordinate of each parity class
in an open rectangle: the count is
\(n^2\operatorname{area}(E)/4+O(n)\). Sum the fixed number of cells
and types. No assumption that \(n\) is even or a multiple of \(G\)
is used.
\end{proof}

The regular grid in the concrete input permits a useful relative
version of the mass estimate.
\begin{lemma}[Mass on every integer grid]\label{lem:massgrid}
For nonnegative four-type fields constant on the open cells of a
\(G\times G\) grid and zero on its grid lines, let
\(A=\tfrac14\sum_r\int F_r\) and
\(A_n=n^{-2}\sum_{p\in V_n}F_{\typ(p)}(p/n)\). For every integer
\(n\ge2G\),
\begin{equation}\label{eq:massgrid}
 A(1-2G/n)^2\le A_n\le A(1+2G/n)^2.
\end{equation}
\end{lemma}
\begin{proof}
For a coordinate parity \(e\in\{0,1\}\), the number of integers of
that parity in \((in/G,(i+1)n/G)\) is
\[
 N_{i,e}=\left\lceil\frac{(i+1)n-eG}{2G}\right\rceil
          -\left\lfloor\frac{in-eG}{2G}\right\rfloor-1.
\]
It lies between \(n/(2G)-1\) and \(n/(2G)+1\). Since the lower
bound is nonnegative, the number of points of any specified parity
type in a cell lies between the squares of these bounds. Multiply
by that cell's nonnegative field value and sum over cells and types.
The continuum contribution of a typed cell of value \(v\) is
\(v/(4G^2)\), giving \eqref{eq:massgrid}. Boundary points have field
zero, as required by the definition of the input.
\end{proof}

\begin{lemma}[Uniform bounded variation]\label{lem:bv}
For a nonnegative field \(F\le M\) on \(R\) fixed rectangular cells,
each of the four functions \(u\mapsto c_{F,j}(x,u)^2\) satisfies
\[
 \|c_{F,j}(x,\cdot)^2\|_\infty\le M^2,
 \qquad \Var(c_{F,j}(x,\cdot)^2)\le(4R+2)M^2,
\]
uniformly in \(x\in[0,1]^2\).
\end{lemma}
\begin{proof}
For \((1,u)\), the chord of a rectangle \((a,b)\times(c,d)\) through
\((x,y)\), for \(u>0\), is
\[
 \pos{\min(b,x+(d-y)/u)-\max(a,x+(c-y)/u)}.
\]
Its changes of branch occur only at slopes from the root to the four
vertices. Commonly refine the at most \(4R\) switch values. On each
interval the total chord has the form \(\alpha+\beta/u\), hence is
monotone, as is its square since it is nonnegative. The square lies
in \([0,M^2]\). At positive switches the values agree. At zero, a line
may coincide with a common boundary; its possible endpoint jump costs
at most one further \(M^2\). The stated bound follows. Reflections
and coordinate exchange give the other three families.
\end{proof}

This proof uses bounded variation, not a uniform derivative bound.
For example, for a half-square and a root a distance \(\epsilon\)
above its boundary, a chord can have the form
\(\max(0,1/2-\epsilon/u)\). Its derivative near \(u=2\epsilon\)
is unbounded as \(\epsilon\downarrow0\), whereas its squared
variation remains at most \(1/4\).

\subsection{A finite sieve for primitive directions}\label{sec:sieve}
Let \(P\) be a fixed finite set of pairwise coprime odd integers greater
than one; primality is unnecessary. Put
\[
 Q=\prod_{p\in P}p,\qquad
 D_P=\frac34\prod_{p\in P}(1-p^{-2}),\qquad
 A_P=\prod_{p\in P}(1+p^{-2}).
\]
For \(e=(\epsilon,\eta)\in\{(0,1),(1,0),(1,1)\}\) and a nonnegative
BV function \(w\) on \([0,1]\), define
\[
 K_{N,e}(w)=\sum_{h\le N}\frac1{h^2}
 \sum_{\substack{1\le a\le h\,;\ (a,h)=1\\(h,a)\bmod2=e}}w(a/h).
\]
Write \(I=\int_0^1w\), \(M_w=\sup w\), and \(V_w=\Var(w)\).

\begin{lemma}[Finite-sieve upper bound]\label{lem:sieve}
For every real \(N\ge Q\),
\begin{equation}\label{eq:sieve}
 K_{N,e}(w)\le\frac{D_P}{3}I\log N+C_P(M_w+V_w),\qquad
 C_P=\left(5+\frac14\log Q\right)A_P.
\end{equation}
The rational upper bound
\(C_P\le[5+\tfrac14\sum_{p\in P}(p-1)]A_P\) is also valid.
\end{lemma}
\begin{proof}
For real \(x\ge1\), omit the coprimality condition and denote the
corresponding sum by \(J_e(x;w)\). Counting alternating tags gives
\[
 \left|\sum_{\substack{1\le b\le k\\b\equiv\eta\ (2)}}w(b/k)
       -\frac k2 I\right|\le2M_w+V_w.
\]
To see this, place intervals of length \(2/k\) to the left of the
tags; the oscillation errors sum to at most \(V_w\), and the omitted
or cut end intervals have total length at most \(4/k\).
Also the parity harmonic sum differs from \(\tfrac12\log x\)
by at most two, by comparing ordinary harmonic sums with integrals.
The finite estimate \(\sum_{k\le m}k^{-2}\le2\) therefore gives
\[
 J_e(x;w)=\frac I4\log x+E_e(x),\qquad
 |E_e(x)|\le5(M_w+V_w).
\]
Nonnegativity of \(w\) allows the bound
\[
 \ind_{(a,h)=1}\le
 \prod_{p\in P}(1-\ind_{p\mid a,\ p\mid h}).
\]
Expand this finite product. For \(A\subseteq P\), put
\(d_A=\prod_{p\in A}p\) and \(s_A=(-1)^{|A|}\).
Pairwise coprimality identifies simultaneous divisibility with
divisibility by \(d_A\); oddness preserves the parity class. Thus
\[
 K_{N,e}(w)\le\sum_{A\subseteq P}\frac{s_A}{d_A^2}J_e(N/d_A;w).
\]
Since \(N\ge Q\), every \(N/d_A\ge1\). Substituting the two-sided
estimate for \(J_e\) and using
\(\sum s_A/d_A^2=4D_P/3\), \(\sum1/d_A^2=A_P\),
\(\log d_A\le\log Q\), and \(I\le M_w\) proves
\eqref{eq:sieve}. It is essential to use a two-sided error estimate
inside the signed inclusion-exclusion sum. Finally \(\log p\le p-1\)
gives the rational constant.
\end{proof}

Only a one-sided bound is needed in this application. Consequently
no infinite Mobius series, value of \(\zeta(2)\), or assertion of
probabilistic independence between parity and direction is required.
The recorded choice of 175 factors gives \(D_P<76/125\); its product
\(Q\) has 440 decimal digits. This large fixed threshold is compatible
with the asymptotic quantifiers, but gives no useful small-grid claim.
The factors and rational comparison are included in the companion data.

\subsection{The common-deadline pair sum}
For four fields \(F\) and a root \(p\), define
\begin{equation}\label{eq:pairsum}
 \mathcal B_n(p;F)=\frac1{2n^2}
 \sum_{\substack{q,z\in V_n\setminus\{p\},\ q\ne z\\p,q,z\text{ collinear}}}
 F_{\typ(q)}(q/n)F_{\typ(z)}(z/n).
\end{equation}
The pair sum is ordered; the factor \(1/2\) removes this duplication.
A superscript \(\Gamma^c\) restricts it to nondesignated directions.

\begin{proposition}[Uniform risk comparison]\label{prop:pair}
If \(D_*\ge D_P\), then for \(n\ge\max(2,Q)\),
\[
 \mathcal B_n^{\Gamma^c}(p;F)
 \le\mathcal R_{\typ(p)}(F;p/n)\log n+C_{\rm geom},
\]
where, with \(E_0=(4R+1)M\), one may take
\begin{equation}\label{eq:geomconstant}
 C_{\rm geom}=M^2+4ME_0+2E_0^2+6C_P(4R+3)M^2.
\end{equation}
This constant is independent of \(n,p\) and of the field values
within their fixed partition and bound.
\end{proposition}
\begin{proof}
For each direction let \(W_v\) be the weighted line sum with the root
removed. The contribution is
\((W_v^2-\sum_qF_{\typ(q)}(q/n)^2)/(2n^2)\).
Discard the nonnegative diagonal term. By \eqref{eq:chorderror},
\(W_v\le nC_v(F_r+F_{r+\bar v})/(2h)+E_0\), while
\(C_v(F_r+F_{r+\bar v})\le2M\). Expanding the square, and using
at most \(4h\) directions at height \(h\le n\), gives
\[
 \mathcal B_n^{\Gamma^c}(p;F)
 \le\frac18\sum_{\substack{v\text{ primitive}/\pm\\h(v)\le n}}
 \frac{C_v(F_r+F_{r+\bar v})(p/n)^2}{h(v)^2}
       +4ME_0+2E_0^2.
\]
Indeed, the cross term is at most
\(\sum_{h\le n}4h\,ME_0/(nh)\le4ME_0\); the squared error is at
most \(\sum_{h\le n}4h\,E_0^2/(2n^2)\le2E_0^2\).
Represent nonaxis directions by the four families with slope
\(a/h\), \(1\le a\le h\). The diagonal overlap can be retained for
an upper bound; the two axes contribute at most \(M^2\) separately.
Apply Lemma~\ref{lem:sieve} to each of the twelve family/parity
functions. Their supremum plus variation is at most
\((4R+3)(2M)^2\), by Lemma~\ref{lem:bv}, so the total sieve error,
including the factor \(1/8\), is at most \(6C_P(4R+3)M^2\).
For each nonzero parity \(v\), the four integrals sum to
\(4g_{F_r+F_{r+v}}\). The coefficient is
\(\tfrac18\cdot\tfrac{D_P}{3}\cdot4=D_P/6\).
Increase it to \(D_*/6\). The directions discarded at the start
only increase the upper bound.
\end{proof}

For a general time-dependent input take \(M=T\) once and for all.
Any smaller bound on its terminal fields is also valid. The same
constant works for every time, including every root's last-use time.
For a fixed root with deadline \(\tau\), finite-product differentiation
gives the exact identity
\begin{equation}\label{eq:timesym}
 \frac1{n^2}\int_0^\tau\sum_{q,z}^{\Gamma^c}m_t(q)F_{\typ(z),t}(z/n)\dd t
 =\mathcal B_n^{\Gamma^c}(p;F_\tau).
\end{equation}
Indeed, differentiate each product \(F_t(q)F_t(z)\), exchange \(q,z\),
and integrate. The index set is symmetric and consists of distinct
points different from the root. Using separate deadlines for \(q\)
and \(z\) would not give this identity.

\subsection{Explicit constants for the 481-phase input}\label{sec:uniform481}
The input defined in Section~\ref{sec:rankinput} has \(G=256\).
Every cumulative field is dominated by its terminal field, whose
maximum is \(6.25798212168<7\). Thus \(R=256^2=65536\) and \(M=7\)
work at every time, although the sum of the phase durations is much
larger than seven. Lemma~\ref{lem:massgrid} gives
\[
 A_n\ge A(1-512/n)^2\qquad(n\ge512).
\]

For the 175 odd primes from 3 through 1049 used by the finite sieve,
\(\sum_{p\in P}(p-1)=84164\). Also
\[
 S:=\sum_{p\in P}p^{-2}<\sum_{k\ge1}\frac1{4k(k+1)}=\frac14.
\]
The elementary symmetric sum of degree \(j\) in the numbers \(p^{-2}\)
is at most \(S^j\); expanding the finite product gives
\(A_P\le(1-S)^{-1}<4/3\). Hence the rational choice
\(C_P=84184/3\) is sufficient in Lemma~\ref{lem:sieve}. Substitution
in \eqref{eq:geomconstant}, with \(E_0=1835015\), gives
\begin{equation}\label{eq:geomnumeric}
 C_{\rm geom}=8897332619623<10^{13}.
\end{equation}

Consequently a continuum certificate
\(\mathcal R_{\typ(p)}(F_{\tau_p};p/n)\le r_*=9999/10000\)
implies, uniformly over every used root,
\begin{equation}\label{eq:discrete481}
 \mathcal B_n^{\Gamma^c}(p;F_{\tau_p})\le a\log n,
 \qquad a=\frac{9999}{10000}+\frac1{320000}
          =\frac{319969}{320000}<\frac{19999}{20000}=d,
\end{equation}
provided \(\log n\ge3\cdot10^{18}\). This condition implies
\(n\ge Q\), since \(\log Q\le84164\), and it absorbs
\eqref{eq:geomnumeric} with room to spare. It is only a sufficient
threshold for the deterministic lattice comparison, not a claimed
threshold for all stochastic error terms in the final theorem.

